\documentclass[a4paper,11pt]{ltjsarticle}
\usepackage[haranoaji]{luatexja-preset}
\usepackage{amsmath,amssymb}
\usepackage[top=12mm,bottom=10mm,left=14mm,right=14mm]{geometry}
\usepackage{array}
\usepackage{tikz}
\usetikzlibrary{calc}
\pagestyle{empty}
\setlength{\parindent}{0pt}
\newlength{\qcol}\setlength{\qcol}{0.47\textwidth}
\newlength{\qgap}
\newlength{\anscolw}\setlength{\anscolw}{0.30\textwidth}
\newlength{\ansh}
\newcommand{\Q}[2]{\parbox[t]{\qcol}{\raggedright (#1)\hspace{0.6em}$\displaystyle #2$}}
\newcommand{\QR}[4]{\Q{#1}{#2}\hfill\Q{#3}{#4}\par\vspace{\qgap}}
\newcommand{\anscell}[1]{\rule[-\ansdrop]{0pt}{\ansh}(#1)}
\newlength{\ansdrop}

\begin{document}
{\large\bfseries 計算問題　23}\hfill 名前(\underline{\hspace{32mm}})\par\vspace{3mm}
■（1）〜（16）の計算をしなさい。（17）〜（21）は方程式を解きなさい。\par\vspace{3mm}
\setlength{\qgap}{5.4mm}
\QR{1}{\dfrac{8}{2}\div3+\dfrac{6}{2}}{2}{(-15)+(-5)-(-8)}
\QR{3}{\dfrac{8}{2}(2x-12)-\dfrac{3}{6}(6x+48)}{4}{16-(-2)\div\left(-\dfrac{2}{3}\right)}
\QR{5}{-8^2+(-2)^2-3}{6}{\dfrac{x-8}{2}\times6}
\QR{7}{(-16)\div(+2)}{8}{(8a+2)+(3a-6)}
\QR{9}{8(2x-2)-(3x-6)}{10}{8\div2\times(-6)}
\QR{11}{\dfrac{1}{8}(16x-24)+\dfrac{1}{6}(18x+12)}{12}{(2x-8)-(6x+3)}
\QR{13}{(16a-6)\div2}{14}{-0.8-(-3.2)-1.6}
\QR{15}{9\times\{2-(-3+1)\times2\}}{16}{(-12)\times(-4)}
\QR{17}{8x-2=-26}{18}{8(x+2)=9x+17}
\QR{19}{0.8x+0.2=-1.4}{20}{\dfrac{1}{4}(x+2)=\dfrac{1}{3}x-1}
\vspace{1mm}
\Q{21}{8(x-2)-2(x+1)=10x-12}\par\vspace{3mm}
\setlength{\ansh}{9.5mm}\setlength{\ansdrop}{6mm}%
\begin{center}\begin{tabular}{|p{\anscolw}|p{\anscolw}|p{\anscolw}|}\hline
\anscell{1} & \anscell{2} & \anscell{3} \\\hline
\anscell{4} & \anscell{5} & \anscell{6} \\\hline
\anscell{7} & \anscell{8} & \anscell{9} \\\hline
\anscell{10} & \anscell{11} & \anscell{12} \\\hline
\anscell{13} & \anscell{14} & \anscell{15} \\\hline
\anscell{16} & \anscell{17} & \anscell{18} \\\hline
\anscell{19} & \anscell{20} & \anscell{21} \\\hline
\end{tabular}\end{center}

\end{document}
